\chapter{Conclusion}

With the industry's huge commitment to the development of \gls{genai} and \gls{llms}, it is important to address the gap between the public adoption of the technology and the public understanding of the technology. This thesis has proposed a taxonomy of misconceptions about \gls{llms}, and has developed a mediation tool aimed at addressing these misconceptions. The results of the study suggest that the mediation tool can be effective in addressing some misconceptions.

As to whether sufficient education has accompanied public adoption, the results suggest that there is still a significant prevalence of misconceptions about \gls{llms} among users. This primarily includes anthropomorphization of the technology, and a lack of understanding of the underlying mechanisms of tokenization, training data and the context window. Some aspects surrounding \gls{llms} are already well communicated, like the environmental impact, but not not to the point of which where it affects use.
The taxonomy proposed in this thesis can be a useful tool for identifying misconceptions for further research and for the development of educational materials.

On addressing the misconceptions, the interactive learning experience provided by the mediation tool developed in this thesis shows promising results in helping users understand some inner workings of \gls{llms}, resolving misconceptions such as anthropomorphization, hallucination blindness and belief of semantic understanding of tokens. However, the gap between the mediation and the complexity of modern \gls{llms} is too large to affect the participants use of the technology in a significant way. 

Future work will involve expanding the mediation tool to bridge this gap, and to further test its effectiveness in addressing misconceptions about \gls{llms}. The taxonomy can also be further refined based on the results of this study, and can be used as a basis for further research on the topic.

As \gls{llms} continue to become further embedded in work, education and daily life, it is important to continue researching the misconceptions surrounding the technology. By identifying and addressing these misconceptions, we can help ensure that users are able to use the technology in a responsible and effective manner. Thus, focus should be shifted from the question of ``Can machines think?'' to the question of ``How do \textit{we} think about machines?''